\begin{helper}[Graph Fiber Sum]
Fix $n\in\mathbb N$, $I\subseteq\mathrm{Fin}(n)$, parameter $M\in\mathbb N$, and $p_0\in\mathbb R$ with $0\le p_0\le1$.
Let $\pi:\mathrm{Fin}(n)\hookrightarrow\mathrm{Fin}(M)\times\mathrm{Fin}(M)$ be an embedding.
Let $P_R=\{\pi(u)_1:u\in I\}$, $P_B=\{\pi(u)_2:u\in I\}$.
Let $\mathcal D$ be the finite set of incident-data families $D=(D_R, D_B)$ where $D_R:\mathrm{Fin}(M)\setminus P_R\to\mathcal{P}(P_R)$ and $D_B:\mathrm{Fin}(M)\setminus P_B\to\mathcal{P}(P_B)$.
For a sample $\omega$ in the cell $\mathcal C_\pi = \{\omega:\noderef{TwoBiteEmbedding}(\omega)=\pi\}$, let its incident data be $D(\omega)$, defined by $D_R(x)=\{r\in P_R:\{x,r\}\text{ is a red edge}\}$ and $D_B(y)=\{b\in P_B:\{y,b\}\text{ is a blue edge}\}$.
Then for any predicate $H$ on the data, we have
\[
\noderef{TwoBiteEventProbability}(n,M,p_0)
 \bigl(\mathcal C_\pi\cap \{\omega:H(D(\omega))\}\bigr)
=
\noderef{TwoBiteEventProbability}(n,M,p_0)(\mathcal C_\pi)
\sum_{\{D\in\mathcal D: H(D)\}} \nu(D),
\]
where $\nu(D)$ is the mass
\[
\prod_{x\notin P_R} p_0^{|D_R(x)|}(1-p_0)^{|P_R|-|D_R(x)|} \prod_{y\notin P_B} p_0^{|D_B(y)|}(1-p_0)^{|P_B|-|D_B(y)|}.
\]
\end{helper}
\begin{proof}
For incident data \(D\), let \(F_D\) be the finite fiber of samples
\(\omega=(G_R,G_B,\pi)\) in \(\mathcal C_\pi\) with \(D(\omega)=D\).  If
\(D\) is not valid, meaning that it is nonempty on \(P_R\) or \(P_B\), or has
some outside value not contained in \(P_R\) or \(P_B\), then no sample in the
cell has this data, so \(F_D=\emptyset\).

By unfolding \(\noderef{TwoBiteEventProbability}\), a sample in the cell has weight
\[
\noderef{GnpGraphWeight}(p_0,G_R)\,
\noderef{GnpGraphWeight}(p_0,G_B)\,
\noderef{UniformInjectionWeight}(n,M),
\]
and the injection factor is constant across the cell.
Partition the red graph edge variables into three classes:
\[
\{r_1,r_2\}\subseteq P_R,\qquad
\{x_1,x_2\}\subseteq \mathrm{Fin}(M)\setminus P_R,\qquad
\{x,r\}\text{ with }x\notin P_R,\ r\in P_R.
\]
The incident class is fixed by \(D_R\): for every \(x\notin P_R\) and
\(r\in P_R\), the edge \(\{x,r\}\) is present exactly when
\(r\in D_R(x)\).  Its contribution to the red
\(\noderef{GnpGraphWeight}\) is therefore
\[
\prod_{x\notin P_R}
p_0^{|D_R(x)|}(1-p_0)^{|P_R|-|D_R(x)|}.
\]

The other two red edge classes are ignored by the incident data.  For each
ignored red unordered pair \(e\), finite summation over its two possible
statuses contributes
\[
(1-p_0)+p_0=1.
\]
Since the ignored red edge set is finite, repeated finite Fubini gives a total
ignored-red contribution of \(1\).  Thus the red graph fiber sum over graphs
with fixed valid red data \(D_R\) is exactly the displayed red incident
factor.  The same partition and calculation for the blue graph gives the blue
incident factor
\[
\prod_{y\notin P_B}
p_0^{|D_B(y)|}(1-p_0)^{|P_B|-|D_B(y)|}.
\]

Multiplying the red factor, the blue factor, and the fixed injection weight,
the total weight of the sample fiber \(F_D\) is
\[
\noderef{UniformInjectionWeight}(n,M)\,\nu(D)
\]
for valid \(D\), and \(0\) for invalid \(D\).

The fibers \(F_D\) partition \(\mathcal C_\pi\).  Summing the previous
identity over all valid \(D\) gives
\[
\noderef{TwoBiteEventProbability}(n,M,p_0)(\mathcal C_\pi)
=
\noderef{UniformInjectionWeight}(n,M)
\sum_{D\in\mathcal D}\nu(D).
\]
The data mass is \(1\): for each outside red vertex \(x\), summing the red
factor over all subsets \(A\subseteq P_R\) gives
\[
\sum_{A\subseteq P_R}p_0^{|A|}(1-p_0)^{|P_R|-|A|}
=(p_0+1-p_0)^{|P_R|}=1,
\]
and the same binomial calculation applies to each outside blue vertex.  Hence
\(\sum_D\nu(D)=1\), and the full cell probability is the fixed injection
weight.

Finally restrict the same fiber partition to the event \(H(D(\omega))\).  The
left side is the sum of the fiber masses over valid data with \(H(D)\), namely
\[
\noderef{UniformInjectionWeight}(n,M)
\sum_{\{D\in\mathcal D:H(D)\}}\nu(D).
\]
Replacing the fixed injection weight by the equal cell probability yields
\[
\noderef{TwoBiteEventProbability}(n,M,p_0)(\mathcal C_\pi)
\sum_{\{D\in\mathcal D:H(D)\}}\nu(D),
\]
which is the claimed formula.
\end{proof}
